% 图 4.5 MNN 架构图（重绘加强版）
% 基于 CICC0903540 技术文档 4.3 - MNN 轻量级推理引擎
\documentclass[tikz,border=10pt]{standalone}

% 强制全局样式与配色
\usepackage{amssymb, amsmath}
\usepackage{fontspec}
\usepackage{xeCJK}
\setCJKmainfont{SimHei}
\usepackage{tikz}
\usetikzlibrary{
  arrows.meta, positioning, shapes.geometric, calc,
  backgrounds, fit, decorations.pathreplacing, patterns
}

% 低饱和度马卡龙配色
\definecolor{mBlue}{HTML}{AEC6CF}
\definecolor{mPink}{HTML}{FFD1DC}
\definecolor{mGreen}{HTML}{B1D8B7}
\definecolor{mPurple}{HTML}{B39EB5}
\definecolor{mYellow}{HTML}{FDFD96}
\definecolor{mGray}{HTML}{E5E4E2}
\definecolor{mDark}{HTML}{4A4A4A}

% 全局 TikZ 样式
\tikzset{
  >=stealth,
  thick,
  draw=mDark,
  box/.style={
    rectangle, rounded corners=3pt, draw=mDark, align=center,
    font=\sffamily\small, minimum height=0.85cm, inner sep=5pt},
  subbox/.style={
    rectangle, rounded corners=2pt, draw=mDark!80, align=center,
    font=\sffamily\scriptsize, minimum height=0.65cm, inner sep=4pt},
  micro/.style={
    rectangle, rounded corners=2pt, draw=mDark!70, align=center,
    font=\sffamily\tiny, minimum height=0.4cm, inner sep=3pt},
  ts/.style={font=\tiny\sffamily, color=mDark!80},
  arr/.style={-{Stealth[length=4pt]}, thick, draw=mDark},
  legendbox/.style={
    rectangle, draw=mDark!50, rounded corners=3pt,
    fill=mGray!15, inner sep=5pt, font=\tiny\sffamily},
}

\begin{document}
\begin{tikzpicture}[
  node distance=0.6cm and 0.8cm,
]
  % ========== 输入模型 ==========
  \node[box, fill=mGray!50, minimum width=12cm] (input)
    {输入模型（TensorFlow / TFLite / Caffe / ONNX）};

  % ========== Converter 模块 ==========
  \node[font=\sffamily\footnotesize\bfseries, color=mDark, below=0.9cm of input.west, anchor=west] (cvtLabel)
    {Converter（模型转换器）};

  \node[subbox, fill=mBlue!40, below=0.3cm of cvtLabel.south west, anchor=north west, minimum width=5.2cm] (frontend)
    {Frontend（前端解析）};

  \node[subbox, fill=mBlue!40, right=0.6cm of frontend, minimum width=5.2cm] (graphopt)
    {Graph Optimizer（图优化器）};

  % Frontend 微观细节
  \node[micro, fill=mBlue!30, below=0.25cm of frontend] (f1) {格式解析};
  \node[micro, fill=mBlue!30, below=0.2cm of f1] (f2) {权重加载};

  % Graph Optimizer 微观细节
  \node[micro, fill=mBlue!30, below=0.25cm of graphopt] (go1) {算子融合};
  \node[micro, fill=mBlue!30, below=0.2cm of go1] (go2) {布局调整};

  % ========== Interpreter 模块 ==========
  \node[font=\sffamily\footnotesize\bfseries, color=mDark, below=1.2cm of cvtLabel.west, anchor=west] (intLabel)
    {Interpreter（推理解释器）};

  \node[subbox, fill=mPurple!40, below=0.3cm of intLabel.south west, anchor=north west, minimum width=5.2cm] (engine)
    {Engine（调度引擎）};

  \node[subbox, fill=mPurple!40, right=0.6cm of engine, minimum width=5.2cm] (backend)
    {Backend（计算后端）};

  % Engine 微观细节
  \node[micro, fill=mPurple!30, below=0.25cm of engine] (en1) {计算图调度};
  \node[micro, fill=mPurple!30, below=0.2cm of en1] (en2) {内存管理};

  % Backend 微观细节
  \node[micro, fill=mPurple!30, below=0.25cm of backend] (be1) {CPU/GPU/NPU};
  \node[micro, fill=mPurple!30, below=0.2cm of be1] (be2) {算子实现};

  % ========== 核心创新：几何计算 + 半自动搜索 ==========
  \node[subbox, fill=mGreen!50, below=0.8cm of engine, minimum width=3.4cm] (geo)
    {几何计算机制};

  \node[subbox, fill=mYellow!50, right=0.5cm of geo, minimum width=3.4cm] (semi)
    {半自动搜索};

  \node[subbox, fill=mPink!40, right=0.5cm of semi, minimum width=2.8cm] (neon)
    {NEON/汇编};

  % 几何计算细节
  \node[micro, fill=mGreen!40, below=0.25cm of geo] (g1) {光栅算子};
  \node[micro, fill=mGreen!40, below=0.2cm of g1] (g2) {region 映射};
  \node[ts, below=0.15cm of g2] {$O(1954) \to O(1055)$};

  % 半自动搜索细节
  \node[micro, fill=mYellow!40, below=0.25cm of semi] (s1) {后端选择};
  \node[micro, fill=mYellow!40, below=0.2cm of s1] (s2) {参数优化};
  \node[ts, below=0.15cm of s2] {tile size, SIMD};

  % NEON/汇编细节
  \node[micro, fill=mPink!30, below=0.25cm of neon] (n1) {ARM NEON};
  \node[micro, fill=mPink!30, below=0.2cm of n1] (n2) {手写汇编};

  % ========== 数据流连线 ==========
  \draw[arr] (input.south) -- ++(0,-0.3) -| (frontend.north);
  \draw[arr] (frontend) -- (graphopt);
  \draw[arr] (graphopt.south) -- ++(0,-0.4) -| (engine.north);
  \draw[arr] (engine) -- (backend);
  \draw[arr] (engine.south) -- ++(0,-0.2) -| (geo.north);
  \draw[arr] (geo) -- (semi);
  \draw[arr] (semi) -- (neon);

  % ========== 层次化表达：背景分组 ==========
  \begin{scope}[on background layer]
    % Converter 区域
    \node[fit=(cvtLabel)(frontend)(graphopt)(f1)(f2)(go1)(go2),
      fill=mBlue!15, draw=mDark!40, dashed,
      rounded corners=5pt, inner sep=10pt] {};

    % Interpreter 区域
    \node[fit=(intLabel)(engine)(backend)(en1)(en2)(be1)(be2)(geo)(semi)(neon)(g1)(g2)(s1)(s2)(n1)(n2),
      fill=mPurple!15, draw=mDark!40, dashed,
      rounded corners=5pt, inner sep=10pt] {};
  \end{scope}

  % ========== 信息密度拉满：算子分类与优化 ==========
  \node[legendbox, anchor=north west] at ([xshift=5pt, yshift=-5pt]current bounding box.north west)
    [align=left, text width=5.0cm] {
    \textbf{MNN 算子分类与优化}\\[2pt]
    • 原子算子（61个）：加、减、乘、除\\
    • 转换算子（45个）：转置、切片、重排\\
    • 复合算子（16个）：卷积、池化、归一化\\
    • 控制流（2个）：if、while\\[2pt]
    \textbf{几何计算收益}：$O(1954) \to O(1055)$\\
    约减少 46\% 后端优化工作量
  };

  % ========== 图例（右下角）==========
  \node[legendbox, anchor=south east] at ([xshift=-5pt, yshift=5pt]current bounding box.south east)
    [align=left, text width=4.8cm] {
    \textbf{图例}\\[1pt]
    \colorbox{mBlue!40}{\rule{0.16cm}{0.12cm}\hspace{0.05cm}} Converter \quad
    \colorbox{mPurple!40}{\rule{0.16cm}{0.12cm}\hspace{0.05cm}} Interpreter\\
    \colorbox{mGreen!50}{\rule{0.16cm}{0.12cm}\hspace{0.05cm}} 几何计算 \quad
    \colorbox{mYellow!50}{\rule{0.16cm}{0.12cm}\hspace{0.05cm}} 半自动搜索\\[1pt]
    Session API: \texttt{resizeTensor} + \texttt{resizeSession}\\
    动态输入支持：$B{\times}3{\times}H{\times}W$（$H,W$ 可变）
  };
\end{tikzpicture}
\end{document}
